\scanpage{25}
(ii) $N'{}^*(X)$ est un \uline{idéal} de $A^*(X)$, donc $\mathrm{Cyc}^*(X)/N^*(X)\simeq A^*(X)/N'{}^*(X)$ est muni d'une structure d'\uline{anneau}, soit $A'{}^*(X)$

(iii) Pour tout hom $f:X\longrightarrow Y$ de \deleted{variétés projectives}\inserted{\unsure{qu. pr.|qu. proj.|q.p.} \deleted{\unsure{préschémas}}} lisses, l'hom $f^*:A^*(Y)-A^*(X)$ applique $N'{}^*(Y)$ dans $N'{}^*(X)$, donc définit par passage au quotient un hom d'anneaux $A'{}^*(Y)-A'{}^*(X)$.

(iv) Pour tout $f$ comme dans (iii), \inserted{avec $f$ \uline{propre},} l'hom \inserted{de groupes} $f_*:A^*(X)\unsure{\longrightarrow|-}A^*(Y)$ \deleted{est} applique $N'{}^*(X)$ dans $N'{}^*(Y)$, et induit donc un hom de groupes (également noté $f_*$)

\[
A'{}^*(X)\longrightarrow A'{}^*(Y)
\]

(v) L'équiv. d'Albanese $\Longrightarrow$ équivalence algébrique. (Donc on peut parler d'\unsure{éléments|élémts} de $A'{}^*(X)$ \uline{alg.} \uline{équiv.} \uline{à} $0$ !)

\uline{Dém}. (i) \deleted{Si $z$ est alb-équiv. \unsure{à} $0$ sur $X$, il est évident qu'il en est de $\unsure{\overline{\mathrm m}|=}$ de $-z$.} \inserted{\unsure{Évident|Évid\textsuperscript{mt}|Évidemment} $0$ est alb-équiv. à $0$.}

Prouvons que si $z,z'$ sont alb-équiv. à $0$ sur $X$, il en est de $\overline{\mathrm m}$ de $z\unsure{-|+}z'$.

Soient $T,T'$ des variétés de paramètres comme dans 1.1. correspondant à $z,z'$. Soit $T_{1}^{\unsure{\prime\prime}}=T\times T'$

% archive sheet number: 24
